\documentclass[10pt,a4paper]{amsart}
\usepackage[margin=19mm]{geometry}
\usepackage{amsmath,amssymb,mathtools}
\usepackage[scr=rsfs,cal=euler]{mathalfa}
\usepackage{xfrac}
\usepackage[unicode,hidelinks,bookmarks=false]{hyperref}
\newcommand{\ie}{i.e.\ }
\newcommand{\cf}{cf.\ }
\newcommand{\resp}{respectively\ }
\newcommand{\Subsection}{\subsection}
\providecommand{\nobreakdash}{\nobreak\hbox{-}}
\setlength{\emergencystretch}{2em}
\multlinegap0pt
\usepackage{amssymb}
\usepackage{tikz}
\usepackage{float}
\newtheorem{theorem}{Theorem}
\title{Factorizations of Matrices With Recursive Entries and Related Topics}
\author{Xiao You Chen, Ali Reza Moghaddamfar, Kambiz Moghaddamfar}
\date{Source: arXiv:2602.07175v1 (CC BY 4.0)}
\begin{document}
\maketitle

\noindent\emph{Source and licence.} Factorizations of Matrices With Recursive Entries and Related Topics. Version arXiv:2602.07175v1 (CC BY 4.0), distributed by arXiv under the Creative Commons Attribution 4.0 International licence, http://creativecommons.org/licenses/by/4.0/, as recorded in the arXiv licence metadata for this exact version and shown on the abstract page https://arxiv.org/abs/2602.07175v1. Full licence notice: \texttt{licenses/arxiv-2602.07175-CC-BY-4.0.txt}.\par\smallskip

\noindent\textit{Editorial scope.} Counted unit: the article's theorem mainth (source0.tex lines 220--230). The article's complete original proof of it is delivered with it (source0.tex lines 231--258, one \texttt{proof} environment of 28 lines). No mathematics is written, repaired or re-derived: the statement and the proof are the article's own lines, in the article's own notation, and no retained line is substituted or re-flowed.  Retained uncounted as the source-local context the proof needs: lines 127--129; lines 198--200.  Nothing was omitted from any retained span, so the package carries no omission record.  No retained line contains a citation, so the wrapper carries no bibliography and no bibliography entry.  No retained line contains a figure environment, an image-inclusion command or drawing code, so no image is delivered and the package declares no asset dependency.  Editorial formatting only, in addition: an AMS \texttt{amsart} preamble that restates the article's own declarations and macros used by the retained lines, namely the environments \texttt{theorem} on separate counters, together with the standard packages those lines need.  The retained environments print in this document as Theorem 1 in order of appearance, whereas the article prints the same items with the numbering of its own full document.  The complete native source of the article as distributed in the arXiv e-print for this exact version is bundled unchanged as \texttt{sources/194156/source0.tex}.
\begin{align}\label{eq1}
P_{\lambda_w, \mu}^{[0, v, w]} \cdot P_{\alpha, \beta}^{[x, y, z]}=P^{[x, vy-wx, vz+w]}_{\eta, \beta},
\end{align}

\begin{align}\label{eq7}
P_{\alpha, \beta}^{[x, y, z]}\cdot (P_{\lambda_w, \mu}^{[0, v, w]})^t=P^{[vx+w, vy-wz, z]}_{\alpha, \eta},
\end{align}

\begin{theorem}[Unifying Factorization]\label{mainth}
For any $r, s, v, w, x, y, z\in \mathbb{C}\) such that $rv\neq 0$, and any two sequences $\alpha, \beta$ with $\alpha_0=\beta_0$, we have the following decomposition:
\[
P_{\alpha, \beta}^{[x, y, z]}=P_{\lambda_s, \mu}^{[0, r, s]}  \cdot P_{\rho, \sigma}^{\left[ \frac{x-w}{v}, \frac{y+xs+zw-sw}{rv}, \frac{z-s}{r} \right]}\cdot (P_{\lambda_w, \mu}^{[0, v, w]} )^t,
\]
where the transformed sequences $\rho=(\rho_i)_{i\geq 0}$ and  $\sigma=(\sigma_i)_{i\geq 0}$ are given by:
\begin{align*}
\rho_{0}=\alpha_0, \ \text{and} \  \rho_{i}= \left(\alpha_i - \sum_{k=0}^{i-1} \binom{i}{k} r^k s^{i-k}  \rho_k\right)/r^i,  \quad \text{for } \  i\geq 1,\\
\sigma_{0}=\beta_0, \ \text{and} \  \sigma_{i}=\left( \beta_i-\sum_{k=0}^{i-1} \binom{i}{k} v^k w^{i-k}  \sigma_k\right)/v^i,  \quad \text{for } \   i\geq 1.
 \end{align*}
\end{theorem}

\begin{proof}
Applying Eq. (\ref{eq1}), we obtain:
\begin{align*}
P_{\lambda_s, \mu}^{[0, r, s]}  \cdot P_{\rho, \sigma}^{\left[ \frac{x-w}{v}, \frac{y+xs+zw-sw}{rv}, \frac{z-s}{r} \right]} =P_{\eta, \sigma}^{[\frac{x-w}{v}, \frac{y+zw}{v}, z]}
\end{align*}
where $\eta=(\eta_i)_{i\geq 0}$ satisfies the following relation:
\begin{align*}
\eta_{i} =\sum_{k=0}^{i} \binom{i}{k} r^k s^{i-k} \rho_{k}, \quad   \text{for}  \   i\geq 0.
\end{align*}
Note that $\eta_i=\alpha_i$ for all $i \geq 0$. In fact, we can demonstrate that:
\begin{align*}
\eta_i&=\sum_{k=0}^{i} \binom{i}{k} r^k s^{i-k} \rho_{k}=r^i \rho_{i}+\sum_{k=0}^{i-1}\binom{i}{k} r^k s^{i-k} \rho_{k}\\ &= \left(\alpha_i-\sum_{k=0}^{i-1} \binom{i}{k} r^k s^{i-k}  \rho_k\right)+\sum_{k=0}^{i-1}\binom{i}{k} r^k s^{i-k} \rho_{k}=\alpha_i.
\end{align*}
Thus, we can express the above equation as:
\begin{align*}
P_{\lambda_s, \mu}^{[0, r, s]}  \cdot P_{\rho, \sigma}^{\left[ \frac{x-w}{v}, \frac{y+xs+zw-sw}{rv}, \frac{z-s}{r} \right]} =P_{\alpha, \sigma}^{[\frac{x-w}{v}, \frac{y+zw}{v}, z]}.
\end{align*}
Applying Eq. (\ref{eq7}) to $P_{\alpha, \sigma}^{[\frac{x-w}{v}, \frac{y+zw}{v}, z]}$, we have:
\begin{align*}
P_{\alpha, \sigma}^{[\frac{x-w}{v}, \frac{y+zw}{v}, z]} \cdot (P_{\lambda_w, \mu}^{[0, v, w]})^t=P_{\alpha, \nu}^{[x, y, z]},
\end{align*}
where $\nu=(\nu_i)_{i\geq 0}$ satisfies the following relation:
\begin{align*}
\nu_i= \sum_{k=0}^{i} \binom{i}{k} v^k w^{i-k} \sigma_{k}, \ \ 
\text{for} \ i\geq 0. 
\end{align*}
 Reasoning as before, we can observe that $\nu_i = \beta_i$ for all $i \geq 0$,  which completes the proof.
\end{proof}

\end{document}
